\documentclass[12pt,a4paper]{article}
\usepackage[utf8]{inputenc}
\usepackage[russian]{babel}
\usepackage{amsmath, amssymb}
\usepackage{geometry}
\geometry{margin=2cm}

\title{Лабораторный журнал моделирования магнитной реснички}
\author{}
\date{}

\begin{document}
\maketitle

\section{Описание модели}

\subsection{Физическая постановка задачи}

В данной работе рассматривается численная модель двухслойной магнитной
реснички, закрепленной на упругой подложке и деформируемой внешним механическим
воздействием. Объект моделирования представляет собой тонкий цилиндрический
стержень малого диаметра \(D\) и полной длины
\[
    L = L_1 + L_2,
\]
состоящий из двух последовательно расположенных по высоте слоев. Нижний слой
длины \(L_1\) выполняется из немагнитного полимера PDMS, а верхний слой длины
\(L_2\) является магнитным композитом на основе PDMS с частицами NdFeB. Радиус
реснички обозначается
\[
    R = \frac{D}{2}.
\]
Ресничка закреплена в цилиндрической подложке радиуса \(R_s\) и толщины
\(H_s\). Подложка служит механическим основанием, через которое задается
закрепление нижней поверхности и передается реакция на внешнее смещение.

Физически нижний PDMS-слой играет роль мягкой эластичной ножки. Он определяет
значительную часть податливости конструкции и позволяет ресничке изгибаться при
поперечном воздействии. Верхний слой содержит остаточную намагниченность,
создающую магнитное поле в окрестности датчика. При деформации реснички меняется
пространственное положение магнитного слоя и, при включенном учете вращения,
локальная ориентация намагниченности. Поэтому магнитное поле в точке или области
измерения изменяется даже при постоянной величине остаточной намагниченности.

Экспериментальная схема интерпретируется как измерение Hall-сенсором магнитного
поля вблизи деформируемой магнитной реснички. До приложения механического
смещения вычисляется поле
\[
    \mathbf B_0 = \mathbf B(\text{начальная конфигурация}),
\]
после деформации вычисляется поле
\[
    \mathbf B_1 = \mathbf B(\text{деформированная конфигурация}).
\]
Основной измеряемой величиной является изменение поля
\[
    \Delta \mathbf B = \mathbf B_1 - \mathbf B_0.
\]
В программных результатах для краткости используются обозначения
\(\mathbf B_0\), \(\mathbf B_1\) и \(d\mathbf B\), где \(d\mathbf B\) совпадает с
\(\Delta \mathbf B\). Компоненты \(dB_x\), \(dB_y\), \(dB_z\) позволяют сравнивать
модель с конкретной ориентацией Hall-сенсора, а норма
\[
    |\Delta \mathbf B| =
    \sqrt{(\Delta B_x)^2 + (\Delta B_y)^2 + (\Delta B_z)^2}
\]
используется как интегральная мера магнитного отклика.

Геометрическая модель является трехмерной. Ресничка и подложка строятся в
геометрическом ядре Gmsh, после чего дискретизуются тетраэдральной сеткой.
Материальная разметка отделяет подложку, нижний PDMS-слой и верхний магнитный
слой. Такая постановка позволяет решать механическую задачу на одной
согласованной сетке, а затем использовать сохраненную конфигурацию для
магнитного постпроцессинга.

\subsection{Механическая модель}

Механическая часть описывает большие деформации мягкого полимерного тела.
Неизвестной является вектор перемещений
\[
    \mathbf u(\mathbf X),
\]
заданный в лагранжевой, то есть исходной, конфигурации. Точка
\(\mathbf X\) после деформации переходит в точку
\[
    \mathbf x = \mathbf X + \mathbf u(\mathbf X).
\]
Основной кинематической величиной является градиент деформации
\[
    \mathbf F = \mathbf I + \nabla \mathbf u,
\]
где \(\mathbf I\) -- единичный тензор. Изменение локального объема описывается
якобианом
\[
    J = \det \mathbf F.
\]
Для физически допустимой деформации требуется \(J>0\), поскольку отрицательный
или нулевой якобиан означал бы вырождение или инверсию элемента сетки.

Материалы описываются сжимаемой неогуковской моделью. Плотность упругой энергии
имеет вид
\[
    \psi =
    \frac{\mu}{2}(I_c - 3)
    -
    \mu \ln J
    +
    \frac{\lambda}{2}(\ln J)^2,
\]
где
\[
    I_c = \mathrm{tr}(\mathbf C), \qquad
    \mathbf C = \mathbf F^{T}\mathbf F
\]
-- первый инвариант правого тензора Коши--Грина. Коэффициенты Ламе выражаются
через модуль Юнга \(E\) и коэффициент Пуассона \(\nu\):
\[
    \mu = \frac{E}{2(1+\nu)}, \qquad
    \lambda =
    \frac{E\nu}{(1+\nu)(1-2\nu)}.
\]
Первый член энергии отвечает за искажение формы, то есть за сдвиговую часть
деформации. Логарифмические члены зависят от \(J\) и штрафуют изменение объема.
Для PDMS и магнитного композита используется почти несжимаемое поведение:
\[
    \nu \approx 0.49.
\]
Такое значение отражает физическое свойство эластомеров сохранять объем при
больших деформациях, но при этом избегает строгого предела \(\nu=0.5\), при
котором стандартная формулировка становится сингулярной.

Напряженное состояние в лагранжевой постановке задается первой тензорной формой
Пиолы--Кирхгофа:
\[
    \mathbf P = \frac{\partial \psi}{\partial \mathbf F}.
\]
Равновесие без объемных сил записывается как
\[
    \mathrm{Div}\,\mathbf P = \mathbf 0.
\]
В слабой форме требуется найти перемещение \(\mathbf u\), удовлетворяющее
граничным условиям, такое что для всех допустимых вариаций \(\mathbf v\)
выполнено
\[
    \int_{\Omega_0} \mathbf P(\mathbf u) : \nabla \mathbf v\,dX = 0.
\]
Здесь \(\Omega_0\) -- исходная область расчета. Нелинейность этой задачи
обусловлена зависимостью \(\mathbf P\) от \(\mathbf F\), а значит от
\(\nabla\mathbf u\).

Нагрузка задается в режиме displacement control. Это означает, что на выбранной
управляющей границе задается предписанное смещение, а реакция системы
вычисляется как результат решения краевой задачи. Такой режим удобен для
валидации, поскольку геометрическое смещение можно задавать устойчиво даже при
мягкой нелинейной механике. Для повышения устойчивости используется continuation
loading: конечное смещение достигается не сразу, а через последовательность
промежуточных шагов
\[
    \Delta x_k = \frac{k}{N}\Delta x,\qquad k=1,\ldots,N.
\]
На каждом шаге решается нелинейная задача, а решение предыдущего шага служит
начальным приближением для следующего.

Механическая задача решается отдельно от магнитной по двум причинам. Во-первых,
механика является более дорогой нелинейной частью, но ее результат можно
сохранить и многократно использовать для разных магнитных параметров. Во-вторых,
магнитная модель в текущем проекте является постпроцессором: она использует
начальную и деформированную геометрию магнитного слоя, но не оказывает обратного
влияния на механическое равновесие. Поэтому вычислительная схема является
односторонне связанной:
\[
    \text{mechanics} \longrightarrow \text{magnetics}.
\]
Механическая часть считается валидированной в рамках проекта: подобраны
параметры сетки и материала, проверены реакции, объемы материальных областей и
качество деформации. Дальнейшие итерации в журнале в основном относятся к
магнитной части.

\subsection{Магнитная модель}

Магнитное поле рассчитывается двумя независимыми способами. Первый способ --
модель суммы объемных диполей, используемая как baseline. Второй способ --
магнитостатическая FEM-модель через скалярный магнитный потенциал. Наличие двух
подходов важно методически: простая дипольная модель дает контроль порядка
величины и помогает обнаруживать грубые ошибки, а FEM-постановка должна лучше
описывать непрерывную магнитостатическую задачу при корректной численной
стабилизации.

\subsubsection{Dipole baseline}

В дипольной модели верхний магнитный слой разбивается на тетраэдральные ячейки
механической сетки. Каждая ячейка \(K\) с объемом \(V_K\) трактуется как малый
магнитный диполь с магнитным моментом
\[
    \mathbf m_K = \mathbf M_K V_K,
\]
где \(\mathbf M_K\) -- намагниченность в ячейке. Если остаточная магнитная
индукция материала равна \(B_r\), то характерная величина намагниченности
оценивается как
\[
    |\mathbf M| \approx \frac{B_r}{\mu_0}.
\]
В начальной конфигурации направление \(\mathbf M\) задается в материальной
системе координат. В деформированной конфигурации возможны два режима: оставить
намагниченность фиксированной в лабораторной системе либо повернуть ее вместе с
локальным вращением материала. Во втором случае из градиента деформации
\(\mathbf F\) выделяется вращательная часть, и намагниченность преобразуется
как материальный вектор.

Поле одного диполя в точке наблюдения задается формулой
\[
    \mathbf B(\mathbf r)
    =
    \frac{\mu_0}{4\pi}
    \left[
    \frac{3\mathbf R(\mathbf m\cdot\mathbf R)}{|\mathbf R|^5}
    -
    \frac{\mathbf m}{|\mathbf R|^3}
    \right],
\]
где
\[
    \mathbf R = \mathbf r - \mathbf r_K
\]
-- вектор от центра диполя к точке наблюдения. Полное поле получается суммой по
всем магнитным ячейкам:
\[
    \mathbf B(\mathbf r) = \sum_{K\subset\Omega_m} \mathbf B_K(\mathbf r).
\]
После решения механики эта сумма вычисляется дважды: для исходных центров и
моментов ячеек, что дает \(\mathbf B_0\), и для деформированных центров и
повернутых моментов, что дает \(\mathbf B_1\).

Достоинство дипольного baseline состоит в простоте и прозрачности физического
смысла. Он не требует построения отдельной воздушной сетки, не содержит
искусственной внешней границы и позволяет быстро проверить порядок величины
магнитного сигнала. Ограничение модели связано с near-field областью. Вблизи
магнитного тела поле сильно меняется на масштабе размера ячейки, а приближение
распределенной намагниченности точечными диполями становится чувствительным к
сетке и расстоянию до датчика. Поэтому baseline не заменяет полноценную
магнитостатическую постановку, но служит важным численным ориентиром.

\subsubsection{Magnetostatic FEM}

В магнитостатической FEM-модели предполагается отсутствие свободных токов в
расчетной области. Тогда уравнения магнитостатики имеют вид
\[
    \nabla\times\mathbf H = 0,
\]
\[
    \mathbf H = -\nabla\varphi,
\]
\[
    \nabla\cdot\mathbf B = 0.
\]
Связь между магнитной индукцией, магнитным полем и намагниченностью задается
соотношением
\[
    \mathbf B = \mu_0(\mathbf H + \mathbf M).
\]
Подставляя \(\mathbf H=-\nabla\varphi\), получаем
\[
    \mathbf B = \mu_0(-\nabla\varphi + \mathbf M),
\]
а из условия \(\nabla\cdot\mathbf B=0\) следует уравнение для скалярного
потенциала:
\[
    \nabla\cdot(-\nabla\varphi + \mathbf M)=0.
\]

Расчетная область \(\Omega\) представляет собой воздушный параллелепипед
\emph{air box}, содержащий начальное и деформированное положение магнитного
слоя, а также область расположения датчика. Намагниченность ненулевая только в
подобласти \(\Omega_m\). Умножая уравнение на тестовую функцию \(v\) и интегрируя
по частям, получаем слабую форму:
\[
    \int_\Omega
    \nabla\varphi\cdot\nabla v\,dx
    =
    \int_{\Omega_m}
    \mathbf M\cdot\nabla v\,dx.
\]
После нахождения \(\varphi\) поле в воздухе вычисляется как
\[
    \mathbf B_{\mathrm{air}} = -\mu_0\nabla\varphi.
\]

Главная численная особенность такой постановки состоит в том, что реальная
магнитная задача является задачей во внешней бесконечной области, тогда как FEM
решается в конечном air box. Поэтому внешняя граница \(\partial\Omega\) является
искусственной. Ранний вариант модели использовал условие
\[
    \varphi|_{\partial\Omega}=0.
\]
Это условие стабилизирует линейную систему, но может заметно искажать поле,
если граница расположена недостаточно далеко от источника. Такая чувствительность
типична для open-boundary magnetic problems: поле убывает в пространстве, но
не обращается точно в ноль на конечном расстоянии, а искусственная граница
задает дополнительное физическое предположение.

В текущих итерациях добавлен режим естественной границы. В этом режиме
потенциал на внешней поверхности не фиксируется. Важно подчеркнуть, что это не
является точным условием бесконечной магнитной области. Это только менее
жесткая аппроксимация open boundary, уменьшающая искусственное заземление
потенциала на конечной границе. Слабая форма тогда соответствует нулевому
естественному потоку на границе, что все равно зависит от размера и формы air
box. Поэтому проверка сходимости по размерам расчетной воздушной области
остается обязательной. При отсутствии условия Дирихле потенциал определен с
точностью до аддитивной константы:
\[
    \varphi \quad\text{и}\quad \varphi + C
\]
дают одинаковое поле \(-\nabla\varphi\). Для устранения этой нулевой моды
используется gauge fixing:
\[
    \varphi(x_{\mathrm{ref}})=0.
\]
Это численное закрепление не задает нулевой потенциал на всей внешней границе,
а лишь выбирает уровень отсчета потенциала.

Дополнительным источником численной чувствительности является модель датчика.
Математическая точка может попасть в отдельную ячейку с локально завышенным или
заниженным градиентом потенциала. Реальный Hall-сенсор имеет конечную
чувствительную область, поэтому более физичной величиной является среднее
\[
    \mathbf B_{\mathrm{sensor}}
    =
    \frac{1}{|\Gamma_s|}
    \int_{\Gamma_s}\mathbf B(\mathbf x)\,dS.
\]
В реализации это среднее аппроксимируется конечным набором точек в плоскости
датчика. Такой sensor averaging уменьшает зависимость результата от одной
ячейки и лучше соответствует экспериментальной интерпретации сигнала.

Еще одна важная часть FEM-модели -- перенос магнитного источника на воздушную
сетку. Механическая сетка и air mesh не являются согласованными. Поэтому
текущая реализация использует cell-center tagging: ячейка воздушной сетки
считается магнитной, если ее центр попадает внутрь одного из тетраэдров
магнитного слоя. Такой подход прост, но может ошибаться в эффективном объеме
источника. Для контроля вводится объемная диагностика:
\[
    V_{\mathrm{source}} =
    \sum_{K\in\Omega_h^m}|K|,
    \qquad
    \varepsilon_V =
    \frac{|V_{\mathrm{source}}-V_{\mathrm{ref}}|}{V_{\mathrm{ref}}}.
\]
Большая ошибка \(\varepsilon_V\) указывает, что поле может быть неверным не из-за
физической модели, а из-за грубой проекции источника на air mesh.

\subsection{Архитектура программной реализации}

Проект организован как набор специализированных модулей, связанных общим
pipeline. Центральная идея архитектуры состоит в разделении механического
решателя, магнитных постпроцессоров и сценариев запуска. Это позволяет
перезапускать магнитные расчеты по сохраненному механическому состоянию без
повторного решения нелинейной механики.

Файл \texttt{main.py} является точкой входа. Он разбирает аргументы командной
строки, настраивает логирование, проверяет доступность runtime DOLFINx/MPI и
передает управление соответствующей функции pipeline в зависимости от режима
\texttt{--mode}. Основные режимы включают механику, полный расчет, дипольный
магнитный постпроцессинг, одиночный magnetostatic FEM и FEM validation sweep.

Файл \texttt{params.py} содержит структуру параметров \texttt{ModelParams},
разбор CLI и загрузку YAML-конфигураций. В этой структуре собраны геометрические
параметры, материальные константы, параметры сетки, настройки механики и
магнитики. Там же определены производные геометрические величины, например
радиус реснички \(R=D/2\) и полная длина \(L=L_1+L_2\). Этот модуль задает
единый интерфейс параметров для остальных частей проекта.

Файл \texttt{mechanics\_model.py} отвечает за построение механической модели.
Он содержит функции создания геометрии Gmsh, генерации тетраэдральной сетки,
назначения материалов, сборки гиперупругой задачи, решения нелинейной системы,
вычисления диагностик и сохранения полей. Именно в этом модуле реализуется
неогуковская механика больших деформаций.

Файл \texttt{magnetics\_dipoles.py} реализует baseline через сумму объемных
диполей. Он извлекает магнитные ячейки из механической сетки, вычисляет их
объемы и центры, оценивает магнитные моменты, учитывает деформацию и при
необходимости поворот намагниченности. Результатом являются поля
\(\mathbf B_0\), \(\mathbf B_1\), \(\Delta\mathbf B\) и сопутствующие
диагностики.

Файл \texttt{magnetics\_fem.py} реализует магнитостатический FEM-постпроцессор.
Он строит air box, переносит магнитный источник на воздушную сетку, решает
задачу для скалярного магнитного потенциала, вычисляет поле в датчике и
формирует расширенную диагностику. В этом модуле находятся численные параметры,
связанные с внешней границей, gauge fixing, sensor averaging, ограничением
размера air mesh и контролем объема источника.

Файл \texttt{pipeline.py} связывает все компоненты в вычислительные сценарии.
Он содержит функции одиночного механического запуска, полного расчета,
магнитного постпроцессинга из restart, validation sweep, записи CSV-summary и
печати кратких отчетов. Этот модуль не задает новую физику, а управляет
последовательностью вычислений и обменом данными между модулями.

\subsection{Pipeline вычислений}

Полный вычислительный процесс можно представить как последовательность этапов.
Сначала по параметрам \(D\), \(L_1\), \(L_2\), \(R_s\), \(H_s\) строится
геометрия в Gmsh. Создаются цилиндрическая подложка, нижний PDMS-слой и верхний
магнитный слой. Затем геометрия дискретизуется тетраэдральной сеткой. Каждой
ячейке присваивается материал: подложка, нижний слой или верхний магнитный
композит.

На следующем этапе решается нелинейная механическая задача. Задается
перемещение управляющей границы, используется continuation loading, и для
каждого шага находится равновесное поле перемещений. После завершения
механического расчета вычисляются реакции, объемы материальных областей,
показатели качества деформации и другие диагностические величины.

Для повторного использования результата сохраняется restart. Он состоит из
трех основных файлов:
\[
\begin{aligned}
&\texttt{mechanics\_restart.npz},\\
&\texttt{params.json},\\
&\texttt{mechanics\_result.json}.
\end{aligned}
\]
Файл \texttt{mechanics\_restart.npz} содержит координаты узлов, тетраэдральные
ячейки, разметку материалов и перемещения в вершинах. Файл \texttt{params.json}
сохраняет параметры модели, а \texttt{mechanics\_result.json} -- механические
диагностики, включая реакцию и объемы слоев.

Магнитный постпроцессинг загружает restart и строит две магнитные конфигурации:
начальную и деформированную. Для каждой из них вычисляется поле в области
датчика:
\[
    \mathbf B_0,\qquad \mathbf B_1.
\]
Затем рассчитывается разность
\[
    \Delta\mathbf B =
    \mathbf B_1-\mathbf B_0.
\]
Если известна механическая реакция \(F\), вводится чувствительность
\[
    S =
    \frac{|\Delta\mathbf B|}{F}.
\]
В коде дополнительно сохраняются компонентные чувствительности, например
\[
    S_x = \frac{\Delta B_x}{F},\qquad
    S_z = \frac{\Delta B_z}{F}.
\]
Результаты записываются в CSV-файлы, что позволяет сравнивать разные параметры
сеток, датчиков, air box и магнитной модели.

\subsection{Численные ограничения}

Модель разрабатывается с учетом ограничений типичного ноутбука с оперативной
памятью порядка \(8\,\mathrm{GB}\). Это ограничение существенно для
магнитостатической FEM-задачи, поскольку air box может занимать объем,
существенно больший, чем сама ресничка. При уменьшении шага \(h_{\mathrm{air}}\)
и увеличении размеров воздушной области число тетраэдров быстро растет, а
вместе с ним растут размер матрицы, стоимость сборки, потребление памяти и
время решения линейной системы.

Поэтому в проекте вводятся laptop-safe параметры. Число ячеек air mesh
ограничивается сверху, а слишком большая сетка должна останавливать расчет до
начала решения. Такой подход важен не только практически, но и методически:
если модель требует гигантского air box для устойчивого результата, проблема
должна решаться постановкой граничных условий, проекцией источника или
постпроцессингом датчика, а не простым неконтролируемым увеличением сетки.

Разделение mechanics и magnetics также является частью численной стратегии.
Нелинейная механика дорогая, но после валидации ее результат можно считать
фиксированным. Магнитные параметры, положение датчика, размер air box,
граничные условия и способ усреднения можно исследовать независимо, загружая
один и тот же restart. Это уменьшает стоимость validation sweep и делает
исследование воспроизводимым.

\subsection{Текущее состояние модели}

На текущем этапе механическая часть считается валидированной. Проверены
геометрия, материальная разметка, реакция на заданное смещение, объемы слоев и
качество деформации. Дипольный baseline реализован и используется как контроль
порядка величины магнитного сигнала.

Магнитостатическая FEM-модель находится на стадии численной стабилизации.
Текущие расчеты показывают, что FEM-сигнал недооценивает экспериментальный
уровень, а результат чувствителен к размеру air box и способу задания внешней
границы. Поэтому основное внимание текущих итераций уделяется четырем
направлениям:
\[
\begin{aligned}
&\text{влияние искусственной магнитной границы},\\
&\text{усреднение поля по конечной области датчика},\\
&\text{качество переноса магнитного источника на air mesh},\\
&\text{сходимость по } h_{\mathrm{air}} \text{ и размерам air box}.
\end{aligned}
\]
Цель этих итераций -- отделить численные ошибки от физических предположений
модели и получить устойчивую основу для дальнейшего сопоставления с
экспериментальным Hall-сигналом.

\section{Итерация: стабилизация магнитостатической FEM-модели}

\subsection*{Обнаруженная проблема}

Механическая часть модели уже прошла валидацию, однако магнитостатическая FEM-модель
пока не является физически валидированной. При \(B_r=0.15\,\mathrm{T}\) дипольный
baseline дает максимум порядка \(|\Delta B|\approx 2.8\,\mu\mathrm{T}\), а текущая
FEM-модель дает еще меньший сигнал, порядка \(|\Delta B|\approx 1.1\,\mu\mathrm{T}\).
Экспериментальная оценка для реакции \(F\approx 60\,\mu\mathrm{N}\) равна
\[
    \Delta B_{\mathrm{target}} \approx 0.63 \cdot 60
    \approx 37.8\,\mu\mathrm{T}.
\]
Следовательно, расхождение не может считаться только малой численной погрешностью.

Исходная магнитостатическая постановка записывается как
\[
    \nabla \times \mathbf H = 0,\qquad
    \mathbf H = -\nabla \varphi,
\]
\[
    \nabla \cdot \mathbf B = 0,\qquad
    \mathbf B = \mu_0(-\nabla\varphi + \mathbf M).
\]
Для скалярного потенциала получается слабая форма
\[
    \int_\Omega \nabla\varphi \cdot \nabla v\,dx
    =
    \int_{\Omega_m} \mathbf M\cdot \nabla v\,dx,
\]
где \(\Omega\) -- расчетная воздушная область, а \(\Omega_m\) -- магнитный слой.

\subsection*{Гипотеза}

Сильная зависимость результата от размеров air box указывает, что условие
\(\varphi=0\) на всей внешней границе слишком жестко приближает бесконечную
магнитную область. Кроме того, поле берется в одной математической точке, что
делает результат чувствительным к локальной ячейке сетки. Наконец, магнитный
источник переносится на воздушную сетку грубым правилом: ячейка считается
магнитной, если ее центр попал внутрь тетраэдра механической магнитной области.
Такой cell-center tagging может давать неверный эффективный объем источника.

\subsection*{Изменения модели и алгоритма}

Добавлен режим внешней границы \texttt{magnetic\_boundary}. Старый режим
\texttt{dirichlet\_zero} сохраняет условие
\[
    \varphi|_{\partial\Omega}=0.
\]
Новый режим \texttt{natural} не задает потенциал на искусственной внешней
границе. В слабой форме это соответствует естественному граничному условию,
которое менее жестко связывает решение с выбранным размером air box. При этом
потенциал \(\varphi\) определен с точностью до константы, поскольку поле зависит
от градиента:
\[
    \mathbf B_{\mathrm{air}} = -\mu_0 \nabla\varphi.
\]
Поэтому добавлена фиксация gauge:
\[
    \varphi(x_{\mathrm{ref}})=0.
\]
В коде используется один закрепленный dof. Это убирает нулевую моду, но не
навязывает нулевой потенциал на всей внешней поверхности.

Добавлено усреднение датчика. Реальный Hall-сенсор измеряет поле не в точке, а
по конечной чувствительной области. Модельная величина записывается как
\[
    \mathbf B_{\mathrm{sensor}}
    =
    \frac{1}{|\Gamma_s|}
    \int_{\Gamma_s} \mathbf B(\mathbf x)\,dS.
\]
В коде интеграл заменен квадратурным средним по набору точек в плоскости
\(z=z_s\). Точки строятся на сетке \(n\times n\) внутри диска радиуса
\(r_s\), после чего вычисляется среднее значение \(\mathbf B\).

Добавлена диагностика переноса магнитного источника. Для начального и
деформированного состояний считается объем ячеек воздушной сетки, помеченных
как магнитные:
\[
    V_{\mathrm{source}} = \sum_{K\in\Omega_h^m} |K|.
\]
Ошибка относительно эталонного объема магнитного слоя оценивается как
\[
    \varepsilon_V =
    \frac{|V_{\mathrm{source}} - V_{\mathrm{ref}}|}{V_{\mathrm{ref}}}.
\]
В качестве \(V_{\mathrm{ref}}\) используется \texttt{volume\_upper\_m3} из
механического результата, а если оно недоступно, аналитическая оценка
\(\pi R^2 L_2\).

Добавлен laptop-safe лимит на размер воздушной сетки. После генерации air mesh
проверяется число тетраэдров. Если оно больше \texttt{max\_air\_cells}, запуск
останавливается до решения линейной системы, если явно не задан
\texttt{allow\_large\_air\_mesh}. Это защищает ноутбук с \(8\,\mathrm{GB}\) RAM
от случайного запуска слишком тяжелой задачи.

\subsection*{Новые параметры}

\[
\begin{aligned}
&\texttt{magnetic\_boundary}\in\{\texttt{dirichlet\_zero},\texttt{natural}\},\\
&\texttt{sensor\_average}\in\{\texttt{false},\texttt{true}\},\\
&\texttt{sensor\_average\_radius},\quad
\texttt{sensor\_average\_n},\\
&\texttt{max\_air\_cells},\quad
\texttt{allow\_large\_air\_mesh}.
\end{aligned}
\]

\subsection*{Проверки}

Нужно проверить одиночный запуск при
\(\texttt{magnetic\_boundary}=\texttt{natural}\), \(h_{\mathrm{air}}=100\,\mu
\mathrm{m}\), малом air box и включенном усреднении датчика. Затем нужно
запустить validation config с безопасными для ноутбука сетками:
\[
    h_{\mathrm{air}} \in \{100,75,50\}\,\mu\mathrm{m},
    \qquad
    R_{\mathrm{air}}/R \in \{8,12,16\}.
\]
В CSV следует смотреть \(\Delta B_x\), \(|\Delta B|\), число air cells,
число использованных точек датчика и ошибки объема источника
\(\varepsilon_V\). Если \(\varepsilon_V>20\%\), перенос магнитного слоя на
air mesh следует считать подозрительным.

\subsection*{Критерии успеха}

Итерация считается успешной, если старые режимы запуска не изменены, FEM
работает как с point sensor, так и с sensor averaging, режим \texttt{natural}
устраняет нулевую моду через gauge, слишком большая air mesh останавливается
понятной ошибкой до solve, а summary содержит диагностику объема магнитного
источника. Физически ожидается, что менее жесткая внешняя граница и усреднение
датчика уменьшат численную нестабильность; достижение экспериментального уровня
сигнала на этой итерации не является обязательным критерием.

\section{Итерация: анализ laptop-safe FEM validation}

\subsection*{Источник данных}

Проанализирован файл
\[
    \texttt{magnetic\_fem\_validation\_results/magnetic\_fem\_validation\_summary.csv}.
\]
Валидационный набор содержит десять расчетов: четыре расчета с изменением
\(h_{\mathrm{air}}\), три расчета с изменением бокового размера air box и три
расчета с изменением нижней границы air box. Все расчеты выполнены при
\[
    B_r = 0.15\,\mathrm{T},\qquad
    x_s/R = 1,\qquad
    z_s = 0,
\]
с естественным граничным условием \texttt{natural}, gauge fixing и усреднением
датчика по 13 точкам в диске радиуса \(25\,\mu\mathrm{m}\).

Целевой экспериментальный масштаб для реакции около \(60\,\mu\mathrm{N}\)
оценивается как
\[
    \Delta B_{\mathrm{target}}
    \approx 0.63\cdot 60
    \approx 37.8\,\mu\mathrm{T}.
\]
В расчетах использованная реакция равна
\[
    F = 59.615\,\mu\mathrm{N}.
\]

\subsection*{Сводка результатов}

Основные численные результаты приведены в таблице. Величина
\(\varepsilon_V\) обозначает максимальную из ошибок объема источника в начальной
и деформированной конфигурациях.

\[
\begin{array}{lrrrrr}
\text{case} &
h_{\mathrm{air}},\,\mu\mathrm{m} &
N_{\mathrm{tet}} &
\max\varepsilon_V,\,\% &
\Delta B_x,\,\mu\mathrm{T} &
|\Delta\mathbf B|,\,\mu\mathrm{T}
\\
\hline
\texttt{h\_air\_150um} & 150 & 11499  & 43.99 & 4.919 & 4.930\\
\texttt{h\_air\_100um} & 100 & 38306  & 6.30  & 5.569 & 5.582\\
\texttt{h\_air\_75um}  & 75  & 88934  & 4.00  & 6.174 & 6.194\\
\texttt{h\_air\_50um}  & 50  & 294762 & 1.41  & 6.071 & 6.077\\
\texttt{air\_radius\_8R}  & 50 & 294762 & 1.41 & 6.071 & 6.077\\
\texttt{air\_radius\_12R} & 50 & 342129 & 1.49 & 7.017 & 7.029\\
\texttt{air\_radius\_16R} & 50 & 605222 & 1.96 & 8.908 & 8.913\\
\texttt{air\_below\_4R}  & 50 & 605222 & 1.96 & 8.908 & 8.913\\
\texttt{air\_below\_8R}  & 50 & 639166 & 1.87 & 7.365 & 7.375\\
\texttt{air\_below\_12R} & 50 & 673902 & 0.80 & 6.768 & 6.784
\end{array}
\]

Максимальный полученный сигнал равен
\[
    |\Delta\mathbf B|_{\max} = 8.913\,\mu\mathrm{T}
\]
для случая \texttt{air\_radius\_16R}. Соответствующая чувствительность равна
\[
    S_{\max}
    =
    \frac{|\Delta\mathbf B|_{\max}}{F}
    \approx
    0.1495\,\mu\mathrm{T}/\mu\mathrm{N}.
\]
Отношение целевого сигнала к максимальному расчетному составляет
\[
    \frac{37.8}{8.913}\approx 4.24.
\]
Следовательно, текущая FEM-модель после введения natural boundary и sensor
averaging стала давать существенно больший сигнал, чем ранний вариант порядка
\(1\,\mu\mathrm{T}\), однако количественно все еще недооценивает эксперимент.

\subsection*{Проверка физической достоверности}

Первый положительный признак состоит в том, что доминирующей компонентой
является \(\Delta B_x\). Во всех расчетах
\[
    |\Delta B_x| \gg |\Delta B_z|,\qquad |\Delta B_y|\ll |\Delta B_x|.
\]
Это согласуется с предыдущим физическим выводом: наиболее сильный сигнал
ожидается около \(x_s\approx R\), \(z_s\approx 0\), а основной вклад связан с
поперечным смещением магнитного слоя. В этом смысле модель воспроизводит
качественную структуру сигнала.

Второй положительный признак связан со сходимостью по \(h_{\mathrm{air}}\).
При фиксированном малом air box переход от \(75\,\mu\mathrm{m}\) к
\(50\,\mu\mathrm{m}\) меняет норму сигнала только примерно на \(1.9\%\):
\[
    |\Delta\mathbf B|_{75}=6.194\,\mu\mathrm{T},\qquad
    |\Delta\mathbf B|_{50}=6.077\,\mu\mathrm{T}.
\]
Это означает, что при данных размерах air box локальное разрешение воздушной
сетки уже близко к приемлемому. Расчет \(h_{\mathrm{air}}=150\,\mu\mathrm{m}\)
нельзя считать физически надежным, так как ошибка объема источника достигает
\[
    \varepsilon_V \approx 44\%.
\]
Он пригоден только как грубый smoke-test, но не как расчет для количественного
сравнения.

Главный отрицательный признак -- отсутствие сходимости по размеру air box.
При \(h_{\mathrm{air}}=50\,\mu\mathrm{m}\) увеличение бокового радиуса от
\(8R\) до \(16R\) увеличивает сигнал:
\[
    6.077\,\mu\mathrm{T}
    \rightarrow
    8.913\,\mu\mathrm{T}.
\]
Относительное изменение составляет порядка \(47\%\) относительно случая \(8R\).
Это слишком много для валидированной магнитостатической модели. Следовательно,
даже natural boundary пока не устраняет полностью влияние искусственной
границы.

Зависимость от нижней границы также остается заметной. При фиксированном
\(R_{\mathrm{air}}=16R\) увеличение \texttt{air\_below\_factor} от 4 до 12
уменьшает сигнал:
\[
    8.913\,\mu\mathrm{T}
    \rightarrow
    6.784\,\mu\mathrm{T}.
\]
Это изменение порядка \(24\%\) относительно \texttt{air\_below\_4R}. Поэтому
нижняя граница также влияет на результат, особенно при расположении датчика в
плоскости \(z_s=0\), близко к основанию реснички.

Диагностика переноса источника показывает, что для расчетов с
\(h_{\mathrm{air}}\le 100\,\mu\mathrm{m}\) ошибка объема в основном находится
ниже \(6.3\%\), а для расчетов \(h_{\mathrm{air}}=50\,\mu\mathrm{m}\) ниже
\(2\%\). Это означает, что в наиболее тяжелых laptop-safe расчетах
cell-center tagging уже передает объем магнитного слоя достаточно хорошо.
Следовательно, текущая недооценка сигнала не объясняется только ошибкой полного
объема источника. Однако корректный объем не гарантирует корректного
распределения источника по поверхности и near-field структуре поля.

\subsection*{Вывод о достоверности}

Текущий результат следует считать физически правдоподобным на качественном
уровне, но не валидированным количественно. Модель правильно воспроизводит
доминирование \(\Delta B_x\), масштаб увеличения сигнала после перехода от
жесткой границы к более мягкой постановке и улучшение проекции источника при
уменьшении \(h_{\mathrm{air}}\). Однако отсутствие сходимости по air box и
расхождение с экспериментальным масштабом более чем в четыре раза показывают,
что magnetostatic FEM пока остается численно чувствительной моделью.

В текущем состоянии результаты можно использовать для сравнительного анализа
параметров внутри одной численной постановки, но нельзя использовать как
окончательную количественную валидацию Hall-сигнала.

\subsection*{Рекомендации с учетом ограничения 8 GB RAM}

Первое практическое решение -- исключить \(h_{\mathrm{air}}=150\,\mu\mathrm{m}\)
из количественного анализа. Этот расчет дешевый, но ошибка объема источника
слишком велика. Для рабочих проверок следует использовать
\[
    h_{\mathrm{air}} = 75\,\mu\mathrm{m}
\]
как быстрый вариант и
\[
    h_{\mathrm{air}} = 50\,\mu\mathrm{m}
\]
как более надежный вариант. При этом число тетраэдров остается ниже текущего
лимита \(700000\) даже для наиболее тяжелого случая \texttt{air\_below\_12R}.

Второе направление -- не увеличивать равномерно air box сверх текущего набора.
Случай \texttt{air\_below\_12R} уже содержит \(673902\) тетраэдра, то есть
практически достигает laptop-safe лимита. Дальнейшее увеличение факторов при
том же равномерном \(h_{\mathrm{air}}=50\,\mu\mathrm{m}\) с высокой вероятностью
приведет к превышению памяти. Вместо этого следует перейти к неравномерной
воздушной сетке: мелкая сетка около магнитного слоя и датчика, более крупная
сетка вдали от источника.

Третье направление -- улучшить аппроксимацию открытой границы без роста числа
ячеек. Возможные варианты:
\[
\begin{aligned}
&\text{сравнение natural и dirichlet\_zero на одинаковых laptop-safe сетках},\\
&\text{добавление более физичного дальнего граничного условия},\\
&\text{использование внешнего грубого буферного слоя},\\
&\text{проверка слабой чувствительности gauge dof к выбору точки закрепления}.
\end{aligned}
\]
Приоритет должен быть у вариантов, которые не требуют миллионов тетраэдров.

Четвертое направление -- улучшить перенос магнитного источника. Несмотря на
малую ошибку полного объема при \(50\,\mu\mathrm{m}\), cell-center tagging
остается грубым near-field приближением. Следующий шаг должен состоять не в
простом уменьшении \(h_{\mathrm{air}}\), а в реализации более точной проекции:
например, volume fraction для ячеек air mesh или локального интегрирования
пересечения магнитного слоя с воздушной ячейкой. Это может дать выигрыш в
физической точности без столь резкого роста памяти.

Пятое направление -- выполнить сопоставление FEM и dipole baseline на одинаковой
модели датчика. Для этого нужно вычислять dipole baseline не только в точке, но
и с тем же sensor averaging. Если расхождение между dipole и FEM сохранится при
одинаковой области усреднения, причина вероятнее всего находится в граничных
условиях FEM или projection source. Если же расхождение уменьшится, существенная
часть различий была связана с моделью датчика.

Шестое направление -- провести малый sweep по радиусу усреднения датчика:
\[
    r_s \in \{10,25,50\}\,\mu\mathrm{m}.
\]
Этот sweep дешевле увеличения air box, потому что использует ту же сетку и
меняет только постпроцессинг. Он позволит понять, насколько измеряемый сигнал
зависит от конечной площади Hall-сенсора.

\subsection*{Критерий следующей успешной итерации}

Следующую итерацию следует считать успешной, если при числе тетраэдров
\[
    N_{\mathrm{tet}} \le 700000
\]
удастся уменьшить зависимость \(|\Delta\mathbf B|\) от размеров air box до
уровня порядка \(10\%\) между соседними laptop-safe случаями, сохранив ошибку
объема источника ниже \(5\%\). Достижение экспериментального значения
\(37.8\,\mu\mathrm{T}\) остается конечной целью, но ближайший численный критерий
должен быть сформулирован как стабилизация результата, а не как подгонка
амплитуды.

\section{Итерация: стабилизация open-boundary и near-field модели}

\subsection*{Уточнение физического смысла граничных условий}

Предыдущие расчеты показали, что переход от \(\varphi=0\) на всей внешней
границе к режиму \texttt{natural} увеличил сигнал и уменьшил искусственное
зажатие потенциала. Однако это не означает, что \texttt{natural} является
физически точной открытой границей. Магнитостатическая задача для локального
намагниченного тела в действительности является задачей во внешней бесконечной
области. В FEM эта область обрезается конечным air box, и на искусственной
границе неизбежно задается приближение.

Условие Дирихле
\[
    \varphi = 0 \quad \text{на } \partial\Omega
\]
можно интерпретировать как грубое ``заземление'' потенциала на конечной границе,
то есть как перенос условия на бесконечности на поверхность air box. Если
граница расположена недостаточно далеко, такое условие может подавлять поле и
уменьшать \(\Delta\mathbf B\).

Естественное условие в слабой форме соответствует нулевому нормальному
градиенту потенциала на искусственной границе:
\[
    \frac{\partial\varphi}{\partial n}=0.
\]
Оно также не является точным условием бесконечной области. Оно лишь менее
жестко ограничивает потенциал, но по-прежнему зависит от размера air box.
Следовательно, оба варианта -- \texttt{dirichlet\_zero} и \texttt{natural} --
являются искусственными truncation boundary conditions. Для магнитостатических
open-boundary problems такая чувствительность к обрезанию области является
стандартной численной трудностью.

Главный вывод текущего анализа состоит в том, что доминирующая неопределенность
сейчас связана не с ошибкой полного объема источника, а с чувствительностью к
границе air domain:
\[
    \varepsilon_V < 2\% \quad \text{в лучших случаях},
\]
но изменение размеров air box все еще меняет \(|\Delta\mathbf B|\) на десятки
процентов.

\subsection*{Graded air mesh}

Равномерная воздушная сетка плохо подходит для laptop-safe расчетов. При
увеличении air box число тетраэдров быстро растет, хотя высокая точность нужна
не во всей области, а вблизи магнитного слоя и датчика. Поэтому введена
неравномерная сетка через background field Gmsh типа \texttt{Box}. Внутри
near-field области используется шаг
\[
    h_{\mathrm{near}},
\]
а вне ее -- более грубый шаг
\[
    h_{\mathrm{far}}.
\]
Near box центрирован вокруг реснички в поперечном направлении:
\[
    |x|\le \alpha R,\qquad |y|\le \alpha R,
\]
где \(\alpha=\texttt{near\_radius\_factor}\). По высоте область refinement
покрывает весь air box:
\[
    -H_{\mathrm{below}}\le z\le L+H_{\mathrm{above}}.
\]
Такой выбор сохраняет мелкое разрешение около магнитного слоя и датчика, но
позволяет уменьшить число дальних тетраэдров.

Добавленные параметры:
\[
\begin{aligned}
&\texttt{h\_air\_near},\\
&\texttt{h\_air\_far},\\
&\texttt{near\_radius\_factor}.
\end{aligned}
\]
По умолчанию
\[
    h_{\mathrm{near}} = h_{\mathrm{air}},\qquad
    h_{\mathrm{far}} = 3h_{\mathrm{air}},\qquad
    \alpha = 3.
\]
Цель этой модификации -- проверять большие air box без перехода к миллионам
элементов. Успехом будет считаться либо уменьшение числа тетраэдров при той же
геометрии air box, либо уменьшение sensitivity к границе при том же laptop-safe
лимите
\[
    N_{\mathrm{tet}}\le 700000.
\]

\subsection*{Согласование FEM и dipole sensor model}

Для корректного сравнения FEM и dipole baseline необходимо использовать
одинаковую модель датчика. Ранее FEM мог использовать sensor averaging, тогда
как dipole baseline в основном вычислялся в одной точке. Это создавало
нечестное сравнение, особенно в near-field области, где поле меняется на
масштабе десятков микрометров.

Введена общая функция построения точек датчика. Для точечного режима
\[
    \Gamma_s = \{\mathbf x_s\},
\]
а для усреднения используется набор точек в диске радиуса \(r_s\) в плоскости
датчика. Тогда и FEM, и dipoles вычисляют одну и ту же величину:
\[
    \mathbf B_{\mathrm{sensor}}
    \approx
    \frac{1}{N_s}\sum_{i=1}^{N_s}\mathbf B(\mathbf x_i).
\]
Это позволит отделить ошибки FEM boundary/source projection от различий,
связанных только с моделью Hall-сенсора.

\subsection*{Sensor averaging sweep}

Добавлен отдельный validation config для проверки зависимости сигнала от
радиуса усреднения:
\[
    r_s\in\{10,25,50,75\}\,\mu\mathrm{m}.
\]
Воздушная область и graded mesh в этом sweep фиксированы. Меняется только
модель датчика. Физически хорошим признаком будет монотонная или объяснимая
зависимость \(\Delta B_x\) и \(|\Delta\mathbf B|\) от \(r_s\), без резких
скачков, не связанных с геометрией поля.

\subsection*{Обновленная стратегия validation}

Расчет \(h_{\mathrm{air}}=150\,\mu\mathrm{m}\) больше не следует использовать
для количественных выводов, так как ранее для него была получена ошибка объема
источника порядка \(44\%\). Основные рабочие сетки текущей итерации:
\[
    h_{\mathrm{near}}=75\,\mu\mathrm{m}
    \quad \text{и} \quad
    h_{\mathrm{near}}=50\,\mu\mathrm{m}.
\]
Главные сравнения по air box должны быть локальными и laptop-safe:
\[
    R_{\mathrm{air}}: 12R \leftrightarrow 16R,
\]
\[
    H_{\mathrm{below}}: 8R \leftrightarrow 12R.
\]
Не следует расширять sweep за пределы \(700000\) тетраэдров. Если для улучшения
open-boundary behavior требуется большая область, предпочтение должно
отдаваться graded mesh и более физичным граничным аппроксимациям, а не
равномерному измельчению сетки.

\subsection*{Критерии успеха}

Следующая итерация считается успешной, если выполнены следующие условия:
\[
\begin{aligned}
&\text{graded mesh уменьшает } N_{\mathrm{tet}}
  \text{ или boundary sensitivity},\\
&\text{FEM и dipoles сравниваются при одинаковом sensor averaging},\\
&\text{sensor averaging sweep дает устойчивую или физически объяснимую зависимость},\\
&\text{изменение }|\Delta\mathbf B|\text{ между лучшими laptop-safe air box}
  \le 10\text{--}15\%,\\
&N_{\mathrm{tet}}\le 700000 \text{ во всех стандартных режимах}.
\end{aligned}
\]

\section{Итерация: анализ результатов graded mesh и sensor averaging}

\subsection*{Источник данных}

Проанализированы три файла результатов:
\[
\begin{aligned}
&\texttt{magnetic\_cilium\_3d\_results\_final/magnetic\_fem\_summary.csv},\\
&\texttt{magnetic\_fem\_validation\_results/magnetic\_fem\_validation\_summary.csv},\\
&\texttt{magnetic\_fem\_sensor\_average\_results/magnetic\_sensor\_average\_validation\_summary.csv}.
\end{aligned}
\]
Все расчеты используют \(\texttt{magnetic\_boundary}=\texttt{natural}\),
sensor averaging и graded air mesh с параметрами \(h_{\mathrm{near}}\),
\(h_{\mathrm{far}}\) и \(\texttt{near\_radius\_factor}=3\).

\subsection*{Основные численные результаты}

Одиночный расчет из \texttt{magnetic\_fem.yaml} дал
\[
    |\Delta\mathbf B| = 4.980\,\mu\mathrm{T},
    \qquad
    \Delta B_x = 4.903\,\mu\mathrm{T},
\]
при очень малом числе ячеек воздушной сетки:
\[
    N_{\mathrm{tet}} = 4146.
\]
Это показывает, что graded mesh действительно резко снижает стоимость расчета.
Однако одновременно возникла большая ошибка проекции деформированного
магнитного источника:
\[
    \varepsilon_{V,0}=2.03\%,\qquad
    \varepsilon_{V,1}=75.08\%.
\]
Число tagged source cells в деформированном состоянии оказалось всего 12 против
94 в начальном состоянии. Такой результат нельзя считать физически надежным:
поле вычисляется для магнитного слоя, объем которого после деформации грубо
искажен на air mesh.

Сводка основных validation cases:
\[
\begin{array}{lrrrrr}
\text{case} &
N_{\mathrm{tet}} &
\varepsilon_{V,0},\% &
\varepsilon_{V,1},\% &
\Delta B_x,\,\mu\mathrm{T} &
|\Delta\mathbf B|,\,\mu\mathrm{T}
\\
\hline
\texttt{h\_air\_100um} & 4146  & 2.03 & 75.08 & 4.903 & 4.980\\
\texttt{h\_air\_75um}  & 9848  & 0.37 & 7.34  & 5.771 & 5.797\\
\texttt{h\_air\_50um}  & 32999 & 0.32 & 4.48  & 6.043 & 6.064\\
\texttt{air\_radius\_12R} & 10906 & 1.90 & 17.96 & 5.275 & 5.379\\
\texttt{air\_radius\_16R} & 14543 & 2.18 & 15.85 & 4.670 & 4.681\\
\texttt{air\_below\_8R}   & 14543 & 2.18 & 15.85 & 4.670 & 4.681\\
\texttt{air\_below\_12R}  & 15562 & 2.26 & 5.41  & 4.906 & 4.915
\end{array}
\]

Максимальный сигнал в новом validation наборе:
\[
    |\Delta\mathbf B|_{\max}=6.064\,\mu\mathrm{T}
\]
для случая \(\texttt{h\_air\_50um}\). Это меньше, чем лучший результат прежней
равномерной сетки \(\approx 8.9\,\mu\mathrm{T}\), но получено при существенно
меньшем числе тетраэдров. Например, для базового air box \(8R/4R/4R\) переход
от равномерной сетки \(h_{\mathrm{air}}=50\,\mu\mathrm{m}\) к graded mesh
уменьшил число ячеек примерно с \(294762\) до \(32999\), то есть почти в 9 раз.

\subsection*{Физическая оценка}

Качественно модель сохраняет правильную структуру сигнала:
\[
    |\Delta B_x| \gg |\Delta B_z|
\]
в большинстве расчетов. Это согласуется с прежним выводом о доминировании
поперечной компоненты около \(x_s/R=1\).

Однако физическая достоверность текущего graded mesh ограничена ошибкой
проекции деформированного источника. Near box задан как
\[
    |x|,|y|\le 3R,
\]
что соответствует только \(180\,\mu\mathrm{m}\) при \(R=60\,\mu\mathrm{m}\).
Деформированный магнитный слой в механическом restart смещается значительно
дальше по \(x\). Поэтому заметная часть деформированного источника оказывается
в области грубой сетки \(h_{\mathrm{far}}\), и cell-center tagging начинает
плохо воспроизводить объем \(\Omega_m\). Это особенно видно в расчетах
\(\texttt{h\_air\_100um}\), \(\texttt{air\_radius\_12R}\) и
\(\texttt{air\_radius\_16R}\), где
\[
    \varepsilon_{V,1}=75.1\%,\quad 18.0\%,\quad 15.9\%.
\]
Следовательно, текущий graded mesh нельзя считать окончательной физической
моделью, несмотря на привлекательную экономию памяти.

Сходимость по air box в новом наборе внешне улучшилась. Для lateral comparison
\[
    |\Delta\mathbf B|:
    5.379\,\mu\mathrm{T}
    \rightarrow
    4.681\,\mu\mathrm{T}
\]
при переходе \(12R\rightarrow 16R\), что соответствует примерно \(13.9\%\)
от среднего значения. Для lower boundary comparison
\[
    4.681\,\mu\mathrm{T}
    \rightarrow
    4.915\,\mu\mathrm{T},
\]
то есть около \(4.9\%\) от среднего значения. Формально это уже близко к
целевому диапазону \(10\text{--}15\%\). Но этот вывод нельзя считать
окончательным, пока ошибка объема деформированного источника превышает
примерно \(5\%\). Иными словами, улучшение air-box sensitivity может быть
частично следствием сглаживания или потери источника на грубой дальней сетке.

\subsection*{Sensor averaging sweep}

В sensor averaging sweep геометрия air box и сетка фиксированы:
\[
    h_{\mathrm{near}}=75\,\mu\mathrm{m},\qquad
    h_{\mathrm{far}}=225\,\mu\mathrm{m},\qquad
    R_{\mathrm{air}}=12R,\qquad
    H_{\mathrm{below}}=8R.
\]
Изменялся только радиус усреднения датчика:
\[
    r_s\in\{10,25,50,75\}\,\mu\mathrm{m}.
\]
Получено:
\[
\begin{array}{lrr}
r_s,\,\mu\mathrm{m} & \Delta B_x,\,\mu\mathrm{T} & |\Delta\mathbf B|,\,\mu\mathrm{T}\\
\hline
10 & 5.279 & 5.384\\
25 & 5.275 & 5.379\\
50 & 5.214 & 5.304\\
75 & 5.117 & 5.195
\end{array}
\]
Зависимость является гладкой и физически объяснимой: увеличение площади датчика
усредняет spatial variation поля и немного уменьшает амплитуду. Относительно
\(r_s=25\,\mu\mathrm{m}\) изменение нормы составляет примерно
\[
    +0.09\%,\quad 0,\quad -1.40\%,\quad -3.42\%.
\]
Следовательно, в текущей геометрии неопределенность от радиуса sensor averaging
мала по сравнению с неопределенностью от air boundary и source projection.

\subsection*{Вывод о достоверности текущей итерации}

Graded mesh как идея подтверждена: она резко уменьшает число тетраэдров и
позволяет запускать большие air box в laptop-safe режиме. Однако конкретная
реализация near box через фиксированный \(\texttt{near\_radius\_factor}=3\)
оказалась недостаточной для деформированной конфигурации. Главная ошибка
текущей итерации -- не boundary condition само по себе, а недостаточное
разрешение области, в которой находится деформированный магнитный слой.

Поэтому результаты текущего graded FEM следует считать полезными для
диагностики численной стратегии, но не окончательными для количественной
валидации Hall-сигнала. Надежными можно считать только те случаи, где
\[
    \varepsilon_{V,1}\lesssim 5\%,
\]
например \(\texttt{h\_air\_50um}\) и частично \(\texttt{air\_below\_12R}\).

\subsection*{Рекомендации для следующей итерации}

Главная рекомендация -- изменить определение near-field области. Она должна
покрывать не только окрестность исходной оси реснички, но и фактическую область
начального и деформированного магнитного слоя, а также sensor averaging disk.
Например, можно задавать эффективный near radius как
\[
    R_{\mathrm{near}}
    =
    \max\left(
    \alpha R,\,
    \max_{\mathbf x\in\Omega_m^0\cup\Omega_m^1}
    \sqrt{x^2+y^2}
    + 2R,\,
    \sqrt{x_s^2+y_s^2}+r_s+R
    \right).
\]
Такой выбор должен сохранить грубую сетку вдали от источника, но не потерять
деформированный магнитный слой.

Вторая рекомендация -- добавить в summary диагностические поля для near mesh:
\[
    R_{\mathrm{near}},\qquad
    h_{\mathrm{near}},\qquad
    h_{\mathrm{far}}.
\]
Это позволит прямо связывать ошибки source projection с тем, попал ли
деформированный магнитный слой в область refinement.

Третья рекомендация -- повторить validation после расширения near box и
сравнить:
\[
\begin{aligned}
&\text{source volume error},\\
&N_{\mathrm{tet}},\\
&|\Delta\mathbf B| \text{ при } 12R\leftrightarrow 16R,\\
&|\Delta\mathbf B| \text{ при } 8R\leftrightarrow 12R \text{ по нижней границе}.
\end{aligned}
\]
Хорошим признаком будет одновременное выполнение условий:
\[
    \varepsilon_{V,1}<5\%,\qquad
    N_{\mathrm{tet}}<700000,\qquad
    \text{air-box sensitivity}<10\text{--}15\%.
\]

Четвертая рекомендация -- оставить sensor averaging radius \(25\,\mu\mathrm{m}\)
как рабочее значение. Sweep показывает, что переход от \(10\) до \(75\,\mu
\mathrm{m}\) меняет \(|\Delta\mathbf B|\) всего на несколько процентов, поэтому
эта неопределенность сейчас не является доминирующей.

Пятая рекомендация -- после исправления near box выполнить честное сравнение
FEM и dipole baseline с одинаковым sensor averaging. Если FEM останется ниже
dipoles при малой ошибке объема источника и слабой air-box sensitivity, тогда
следующим кандидатом на ошибку будет слабая форма, boundary approximation или
способ представления \(\mathbf M\) в air mesh.

\end{document}
